\chapter{Conclusioni}
\label{chap:conclusioni}
\section{Consuntivo finale}
L'attività di stage si è svolta coerentemente con il preventivo definito, per una durata complessiva di 300 ore. Si sono tuttavia verificati alcuni scostamenti rispetto alla pianificazione iniziale delle 
singole attività. In particolare, la realizzazione del motore di alert e delle relative regole, per la quale erano state previste 45 ore, ha richiesto un impegno inferiore rispetto a quanto preventivato. 
Al contrario, l'impegno necessario per l'implementazione del modulo XAI era stato inizialmente sottostimato e, a causa della maggiore complessità delle funzionalità richieste e di alcune problematiche 
relative all'infrastruttura, ha richiesto un numero di ore superiore a quello previsto.
\newline
Nonostante tali variazioni, le differenze tra le singole attività si sono complessivamente compensate, permettendo di mantenere il carico di lavoro totale coerente con quanto previsto dal piano di lavoro e 
di completare lo stage entro le 300 ore stabilite.

\newpage
\section{Raggiungimento degli obiettivi}
\begin{center}
    \keepXColumns
    \rowcolors{2}{lightblue}{white}
        \begin{longtable}{|l|p{0.45\textwidth}|l|l|}
            \hline
            \rowcolor{gold}
            \textbf{Codice} & \textbf{Descrizione requisito} & \textbf{Tipologia} & \textbf{Soddisfatto}\\
            \hline
            \endhead

            \hline
            \rowcolor{white}
            \caption{Tabella riguardante lo stato dei requisiti dell'attività di stage, aggiornati a fine attività di stage} \label{tab:requisiti_stage_finali} \\
            \endlastfoot 

            \textbf{O01} & Implementare un sistema XAI auditabile per registrare e consultare decisioni, fattori, evidenze e snapshot dei dati. & Obbligatorio & Sì\\
            \hline
            \textbf{O02} & Realizzare regole alert, eventi e notifiche configurabili per segnalare proattivamente criticità operative. & Obbligatorio & Sì\\
            \hline
            \textbf{O03} & Integrare i punti decisionali della terza consegna nei log XAI e negli alert, limitandosi a controlli, riprese in mano e integrazioni leggere. & Obbligatorio & Sì\\
            \hline
            \textbf{D01} & Differenziare le spiegazioni XAI e i dettagli visibili in base al ruolo dell'utente. & Desiderabile & Sì\\
            \hline
            \textbf{D02} & Predisporre run di ottimizzazione tracciabili con input, output, score, metriche e log decisionale associato. & Desiderabile & No\\
            \hline
            \textbf{D03} & Introdurre dashboard o riepiloghi sugli alert generati, presi in carico e risolti. & Desiderabile & Sì\\
            \hline
            \textbf{F01} & Collegare un solver o microservizio esterno per scenari specialistici come nesting laser o sequenziamento avanzato. & Funzionale & No\\
            \hline
            \textbf{F02} & Aggiungere suggerimenti automatici di configurazione soglie basati sui dati storici. & Funzionale & No\\
            \hline
            \textbf{F03} & Integrare le spiegazioni XAI nel chatbot sviluppato da un altro piano di lavoro o modulo parallelo. & Funzionale & No\\
            \hline
        \end{longtable}
\end{center}

Come mostrato dalla tabella \ref{tab:requisiti_stage_finali}, la maggior parte dei requisiti individuati è stata soddisfatta, ad eccezione dei requisiti facoltativi. Questi ultimi non sono stati 
implementati principalmente a causa della necessità di approfondire componenti del sistema che non risultavano direttamente coinvolte nella realizzazione dei requisiti obbligatori.
\newline
L'implementazione di tali requisiti avrebbe infatti richiesto un'ulteriore fase di analisi e audit dell'architettura esistente, necessaria per acquisire una conoscenza sufficiente delle componenti 
interessate e individuare le modalità di integrazione più appropriate. Considerata la durata limitata del periodo di stage, si è pertanto scelto di concentrare le attività sui requisiti obbligatori, 
garantendone la completa implementazione e validazione.


\section{Conoscenze acquisite}
L'attività di stage ha rappresentato un'importante esperienza formativa, permettendo di consolidare conoscenze già acquisite durante il percorso di studi e di sviluppare nuove competenze in ambito tecnico e 
metodologico.
\newline
In particolare, sono state approfondite le conoscenze relative alla programmazione orientata agli oggetti, con un focus sul linguaggio PHP e sui framework Laravel e Filament. L'utilizzo di tali tecnologie 
all'interno di un'applicazione reale ha permesso di comprendere più approfonditamente l'organizzazione e l'interazione tra le diverse componenti di un sistema software, consentendo di intervenire direttamente 
sul backend, sull'interfaccia utente e sul livello di persistenza dei dati.
\newline
Un'ulteriore importante opportunità formativa è stata rappresentata dallo sviluppo di funzionalità basate sull'intelligenza artificiale. In particolare, l'implementazione del modulo XAI ha permesso di acquisire 
esperienza nell'integrazione di modelli linguistici all'interno di applicazioni esistenti e nell'utilizzo delle API messe a disposizione dai relativi provider. Questo ha consentito di approfondire un ambito 
tecnologico di particolare attualità e rilevanza nel contesto dello sviluppo software.
\newline
Infine, l'attività ha permesso di acquisire competenze relative alle metodologie e agli strumenti utilizzati nel processo di sviluppo software. In particolare, sono state approfondite le pratiche di \textbf{Continuous 
Integration} e di testing automatico, oltre all'utilizzo sistematico degli \textit{Issue Tracking System} come strumento per la pianificazione, l'organizzazione e il monitoraggio delle attività di sviluppo. È stato 
inoltre possibile sperimentare direttamente l'utilizzo del framework \textbf{Scrum} come metodologia di lavoro, partecipando alla pianificazione e all'organizzazione delle attività all'interno di sprint di sviluppo.

\section{Valutazione personale}
Il periodo di tirocinio ha rappresentato un'importante occasione di crescita personale e professionale, sia dal punto di vista dello sviluppo delle competenze tecniche sia per quanto riguarda l'approccio al lavoro 
e alla collaborazione all'interno di un team. La possibilità di lavorare su un progetto di reale interesse e destinato a un contesto aziendale concreto ha costituito un'importante fonte di motivazione e ha 
contribuito a rendere l'esperienza particolarmente stimolante e soddisfacente.
\newline
L'esperienza presso Spazio Dev S.r.l. è stata inoltre molto positiva dal punto di vista dell'ambiente lavorativo. Le dimensioni contenute dell'azienda hanno rappresentato un'opportunità significativa, permettendomi 
di avere una visione più ampia del processo di sviluppo e di osservare direttamente le diverse attività e responsabilità coinvolte nella realizzazione del progetto. Questo ha contribuito a comprendere meglio non 
solo gli aspetti strettamente legati alla programmazione, ma anche le dinamiche organizzative e collaborative che caratterizzano lo sviluppo di un prodotto software.
\newline
Un ulteriore aspetto particolarmente significativo è stato il rapporto instaurato con gli sviluppatori che mi hanno seguito durante il periodo di stage. Il confronto quotidiano e il supporto ricevuto hanno 
rappresentato una parte importante del percorso formativo, permettendomi di apprendere nuove metodologie di lavoro, confrontarmi con problematiche concrete e migliorare progressivamente le mie competenze sia 
tecniche sia relazionali.




\newpage